% birkon/sections/birkat-hamazon.tex
% Birkat HaMazon: Edot HaMizrach base text, edited toward the Seattle Turkish
% nusach. Source, license, every change from the source, and Yosef's nusach
% decisions are recorded in birkon/sources/birkat-hamazon.md.

\sectiontitle{בִּרְכַּת הַמָּזוֹן}

\rubric{Wash mayim aharonim. It is good to say:}
\textbf{לַמְנַצֵּח} בִּנְגִינֹת מִזְמוֹר שִׁיר: אֱלֹהִים יְחָנֵּנוּ וִיבָרְכֵנוּ יָאֵר פָּנָיו אִתָּנוּ סֶלָה: לָדַעַת בָּאָרֶץ דַּרְכֶּךָ בְּכָל גּוֹיִם יְשׁוּעָתֶךָ: יוֹדוּךָ עַמִּים אֱלֹהִים יוֹדוּךָ עַמִּים כֻּלָּם: יִשְׂמְחוּ וִירַנְּנוּ לְאֻמִּים כִּי תִשְׁפֹּט עַמִּים מִישׁוֹר וּלְאֻמִּים בָּאָרֶץ תַּנְחֵם סֶלָה: יוֹדוּךָ עַמִּים אֱלֹהִים יוֹדוּךָ עַמִּים כֻּלָּם: אֶרֶץ נָתְנָה יְבוּלָהּ יְבָרְכֵנוּ אֱלֹהִים אֱלֹהֵינוּ: יְבָרְכֵנוּ אֱלֹהִים וְיִירְאוּ אֹתוֹ כָּל אַפְסֵי אָרֶץ:

\textbf{אֲבָרְכָה} אֶת \Shem{} בְּכָל עֵת תָּמִיד תְּהִלָּתוֹ בְּפִי: סוֹף דָּבָר הַכֹּל נִשְׁמָע אֶת הָאֱלֹהִים יְרָא וְאֶת מִצְוֹתָיו שְׁמֹר כִּי זֶה כָּל הָאָדָם: תְּהִלַּת \Shem{} יְדַבֵּר פִּי וִיבָרֵךְ כָּל בָּשָׂר שֵׁם קָדְשׁוֹ לְעוֹלָם וָעֶד: וַאֲנַחְנוּ נְבָרֵךְ יָה מֵעַתָּה וְעַד עוֹלָם הַלְלוּיָהּ: (רַגְלִי עָמְדָה בְמִישׁוֹר, בְּמַקְהֵלִים אֲבָרֵךְ \Shem{}):

\divider

\rubric{When three or more have eaten together, they make a zimun.}
\rubric{The leader says:}
הַב לָן וְנִבְרִיךְ לְמַלְכָּא עִלָּאָה קַדִּישָׁא.

\rubric{The others respond:}
שָׁמַיִם.

\rubric{The leader continues:}
בִּרְשׁוּת מַלְכָּא עִלָּאָה קַדִּישָׁא

\begin{seasonal}{On Shabbat add:}
וּבִרְשׁוּת שַׁבָּת מַלְכְּתָא
\end{seasonal}

\begin{seasonal}{On Yom Tov add:}
וּבִרְשׁוּת יוֹמָא טָבָא אֻשְׁפִּיזָא קַדִּישָׁא
\end{seasonal}

\rubric{Some add the words in brackets:}
[וּבִרְשׁוּת מוֹרַי וְרַבּוֹתַי] וּבִרְשׁוּתְכֶם. וַיְדַבֵּר אֵלַי, זֶה הַשֻּׁלְחָן אֲשֶׁר לִפְנֵי \Shem{}:

\rubric{The leader says (with ten or more, add the word in brackets):}
נְבָרֵךְ [אֱלֹהֵינוּ] שֶׁאָכַלְנוּ מִשֶּׁלּוֹ.

\rubric{The others respond. With ten or more, add the first word in brackets; on Shabbat, add the second:}
בָּרוּךְ [אֱלֹהֵינוּ] שֶׁאָכַלְנוּ מִשֶּׁלּוֹ וּבְטוּבוֹ הַגָּדוֹל [וְהַקָּדוֹשׁ] חָיִינוּ.

\rubric{The leader repeats:}
בָּרוּךְ [אֱלֹהֵינוּ] שֶׁאָכַלְנוּ מִשֶּׁלּוֹ וּבְטוּבוֹ הַגָּדוֹל [וְהַקָּדוֹשׁ] חָיִינוּ.

\divider

\textbf{בָּרוּךְ} אַתָּה \Shem{} אֱלֹהֵינוּ מֶלֶךְ הָעוֹלָם, (הַזָּנֵנוּ וְלֹא מִמַּעֲשֵׂנוּ, הַמְפַרְנְסֵנוּ וְלֹא מִצִּדְקוֹתֵינוּ, הַמַּעֲדִיף טוּבוֹ עָלֵינוּ) הַזָּן אֶת הָעוֹלָם כֻּלּוֹ בְּטוּבוֹ, בְּחֵן בְּחֶסֶד וּבְרַחֲמִים הוּא נוֹתֵן לֶחֶם לְכָל בָּשָׂר, כִּי לְעוֹלָם חַסְדּוֹ. וּבְטוּבוֹ הַגָּדוֹל תָּמִיד לֹא חָסַר לָנוּ, וְאַל יֶחְסַר לָנוּ מָזוֹן תָּמִיד לְעוֹלָם וָעֶד, כִּי הוּא אֵל זָן וּמְפַרְנֵס לַכֹּל. וְשֻׁלְחָנוֹ עָרוּךְ לַכֹּל, וְהִתְקִין מִחְיָה וּמָזוֹן לְכָל בְּרִיּוֹתָיו אֲשֶׁר בָּרָא בְּרַחֲמָיו וּבְרוֹב חֲסָדָיו, כָּאָמוּר, \textbf{פּ}וֹתֵחַ \textbf{אֶ}ת \textbf{יָ}דֶךָ, וּמַשְׂבִּיעַ לְכָל חַי רָצוֹן. בָּרוּךְ אַתָּה \Shem{}, הַזָּן בְּרַחֲמָיו אֶת הַכֹּל.

(עַל אַרְצֵנוּ וְעַל נַחֲלַת אֲבוֹתֵינוּ) \textbf{נוֹדֶה} לְךָ \Shem{} אֱלֹהֵינוּ, עַל שֶׁהִנְחַלְתָּ לַאֲבוֹתֵינוּ אֶרֶץ חֶמְדָה טוֹבָה וּרְחָבָה, בְּרִית וְתוֹרָה, חַיִּים וּמָזוֹן. עַל שֶׁהוֹצֵאתָנוּ מֵאֶרֶץ מִצְרַיִם, וּפְדִיתָנוּ מִבֵּית עֲבָדִים. וְעַל בְּרִיתְךָ שֶׁחָתַמְתָּ בִּבְשָׂרֵנוּ. וְעַל תּוֹרָתְךָ שֶׁלִּמַּדְתָּנוּ. וְעַל חֻקֵּי רְצוֹנָךְ שֶׁהוֹדַעְתָּנוּ. וְעַל חַיִּים וּמָזוֹן שֶׁאַתָּה זָן וּמְפַרְנֵס אוֹתָנוּ:

\begin{seasonal}{On Hanukkah and Purim:}
\textbf{עַל הַנִּסִּים} וְעַל הַפֻּרְקָן וְעַל הַגְּבוּרוֹת וְעַל הַתְּשׁוּעוֹת וְעַל הַנִּפְלָאוֹת וְעַל הַנֶּחָמוֹת שֶׁעָשִׂיתָ לַאֲבוֹתֵינוּ בַּיָּמִים הָהֵם בַּזְּמַן הַזֶּה.

\rubric{On Hanukkah:}
\textbf{בִּימֵי מַתִּתְיָה} בֶּן יוֹחָנָן כֹּהֵן גָּדוֹל חַשְׁמוֹנָאִי וּבָנָיו, כְּשֶׁעָמְדָה מַלְכוּת יָוָן הָרְשָׁעָה עַל עַמְּךָ יִשְׂרָאֵל, לְשַׁכְּחָם תּוֹרָתֶךָ, וּלְהַעֲבִירָם מֵחֻקֵּי רְצוֹנָךְ. וְאַתָּה בְּרַחֲמֶיךָ הָרַבִּים, עָמַדְתָּ לָהֶם בְּעֵת צָרָתָם, רַבְתָּ אֶת רִיבָם, דַּנְתָּ אֶת דִּינָם, נָקַמְתָּ אֶת נִקְמָתָם. מָסַרְתָּ גִבּוֹרִים בְּיַד חַלָּשִׁים, וְרַבִּים בְּיַד מְעַטִּים, וּרְשָׁעִים בְּיַד צַדִּיקִים, וּטְמֵאִים בְּיַד טְהוֹרִים, וְזֵדִים בְּיַד עוֹסְקֵי תוֹרָתֶךָ, וּלְךָ עָשִׂיתָ שֵׁם גָּדוֹל וְקָדוֹשׁ בְּעוֹלָמָךְ, וּלְעַמְּךָ יִשְׂרָאֵל עָשִׂיתָ תְּשׁוּעָה גְדוֹלָה וּפֻרְקָן כְּהַיּוֹם הַזֶּה. וְאַחַר כָּךְ בָּאוּ בָנֶיךָ לִדְבִיר בֵּיתֶךָ, וּפִנּוּ אֶת הֵיכָלֶךָ, וְטִהֲרוּ אֶת מִקְדָּשֶׁךָ, וְהִדְלִיקוּ נֵרוֹת בְּחַצְרוֹת קָדְשֶׁךָ, וְקָבְעוּ שְׁמוֹנָה יָמִים אֵלּוּ בְּהַלֵּל וּבְהוֹדָאָה, וְעָשִׂיתָ עִמָּהֶם נִסִּים וְנִפְלָאוֹת, וְנוֹדֶה לְשִׁמְךָ הַגָּדוֹל, סֶלָה.

\rubric{On Purim:}
\textbf{בִּימֵי מָרְדְּכַי} וְאֶסְתֵּר בְּשׁוּשַׁן הַבִּירָה כְּשֶׁעָמַד עֲלֵיהֶם הָמָן הָרָשָׁע בִּקֵּשׁ לְהַשְׁמִיד לַהֲרוֹג וּלְאַבֵּד אֶת כָּל הַיְּהוּדִים מִנַּעַר וְעַד זָקֵן טַף וְנָשִׁים בְּיוֹם אֶחָד בִּשְׁלֹשָׁה עָשָׂר לְחֹדֶשׁ שְׁנֵים עָשָׂר הוּא חֹדֶשׁ אֲדָר וּשְׁלָלָם לָבוֹז וְאַתָּה בְּרַחֲמֶיךָ הָרַבִּים הֵפַרְתָּ אֶת עֲצָתוֹ וְקִלְקַלְתָּ אֶת מַחֲשַׁבְתּוֹ וַהֲשֵׁבוֹתָ לוֹ גְמוּלוֹ בְרֹאשׁוֹ וְתָלוּ אוֹתוֹ וְאֶת בָּנָיו עַל הָעֵץ וְעָשִׂיתָ עִמָּהֶם נִסִּים וְנִפְלָאוֹת וְנוֹדֶה לְשִׁמְךָ הַגָּדוֹל סֶלָה:
\end{seasonal}

עַל הַכֹּל \Shem{} אֱלֹהֵינוּ אֲנַחְנוּ מוֹדִים לָךְ וּמְבָרְכִים אֶת שְׁמָךְ כָּאָמוּר וְאָכַלְתָּ וְשָׂבָעְתָּ. וּבֵרַכְתָּ אֶת \Shem{} אֱלֹהֶיךָ עַל הָאָרֶץ הַטּוֹבָה אֲשֶׁר נָתַן לָךְ: בָּרוּךְ אַתָּה \Shem{} עַל הָאָרֶץ וְעַל הַמָּזוֹן:

\textbf{רַחֵם} \Shem{} אֱלֹהֵינוּ עָלֵינוּ וְעַל יִשְׂרָאֵל עַמָּךְ. וְעַל יְרוּשָׁלַיִם עִירָךְ. וְעַל הַר צִיּוֹן מִשְׁכַּן כְּבוֹדָךְ וְעַל הֵיכָלָךְ. וְעַל מְעוֹנָךְ. וְעַל דְּבִירָךְ. וְעַל הַבַּיִת הַגָּדוֹל וְהַקָּדוֹשׁ שֶׁנִּקְרָא שִׁמְךָ עָלָיו. אָבִינוּ רְעֵנוּ זוּנֵנוּ. פַּרְנְסֵנוּ כַּלְכְּלֵנוּ. הַרְוִיחֵנוּ הַרְוַח לָנוּ מְהֵרָה מִכָּל צָרוֹתֵינוּ. וְנָא אַל תַּצְרִיכֵנוּ \Shem{} אֱלֹהֵינוּ לִידֵי מַתְּנוֹת בָּשָׂר וָדָם. וְלֹא לִידֵי הַלְוָאָתָם. אֶלָּא לְיָדְךָ הַמְּלֵאָה וְהָרְחָבָה. הָעֲשִׁירָה וְהַפְּתוּחָה. יְהִי רָצוֹן שֶׁלֹּא נֵבוֹשׁ בָּעוֹלָם הַזֶּה. וְלֹא נִכָּלֵם לְעוֹלָם הַבָּא. וּמַלְכוּת בֵּית דָּוִד מְשִׁיחָךְ תַּחֲזִירֶנָּה לִמְקוֹמָהּ בִּמְהֵרָה בְיָמֵינוּ:

\begin{seasonal}{On Shabbat:}
\textbf{רְצֵה וְהַחֲלִיצֵנוּ} \Shem{} אֱלֹהֵינוּ בְּמִצְוֹתֶיךָ וּבְמִצְוַת יוֹם הַשְּׁבִיעִי. הַשַּׁבָּת הַגָּדוֹל וְהַקָּדוֹשׁ הַזֶּה כִּי יוֹם גָּדוֹל וְקָדוֹשׁ הוּא מִלְּפָנֶיךָ. נִשְׁבּוֹת בּוֹ וְנָנוּחַ בּוֹ וְנִתְעַנֵּג בּוֹ כְּמִצְוַת חֻקֵּי רְצוֹנָךְ. וְאַל תְּהִי צָרָה וְיָגוֹן בְּיוֹם מְנוּחָתֵנוּ. וְהַרְאֵנוּ בְּנֶחָמַת צִיּוֹן בִּמְהֵרָה בְיָמֵינוּ. כִּי אַתָּה הוּא בַּעַל הַנֶּחָמוֹת וְאַף עַל פִּי שֶׁאָכַלְנוּ וְשָׁתִינוּ חָרְבַּן בֵּיתְךָ הַגָּדוֹל וְהַקָּדוֹשׁ לֹא שָׁכַחְנוּ. אַל תִּשְׁכָּחֵנוּ לָנֶצַח וְאַל תִּזְנָחֵנוּ לָעַד כִּי אֵל מֶלֶךְ גָּדוֹל וְקָדוֹשׁ אָתָּה:
\end{seasonal}

\begin{seasonal}{On Rosh Hodesh, festivals, and Hol HaMoed:}
אֱלֹהֵינוּ וֵאלֹהֵי אֲבוֹתֵינוּ, יַעֲלֶה וְיָבֹא, וְיַגִּיעַ וְיֵרָאֶה וְיֵרָצֶה, יִשָּׁמַע יִפָּקֵד וְיִזָּכֵר, זִכְרוֹנֵנוּ וְזִכְרוֹן אֲבוֹתֵינוּ, זִכְרוֹן יְרוּשָׁלַיִם עִירָךְ, וְזִכְרוֹן מָשִׁיחַ בֶּן דָּוִד עַבְדָּךְ, וְזִכְרוֹן כָּל עַמְּךָ בֵּית יִשְׂרָאֵל לְפָנֶיךָ, לִפְלֵטָה, לְטוֹבָה, לְחֵן, לְחֶסֶד וּלְרַחֲמִים לְחַיִּים טוֹבִים וּלְשָׁלוֹם בְּיוֹם:

\rubric{On Rosh Hodesh:}
רֹאשׁ חֹדֶשׁ הַזֶּה,

\rubric{On Pesah:}
חַג הַמַּצּוֹת הַזֶּה, בְּיוֹם [טוֹב] מִקְרָא קֹדֶשׁ הַזֶּה,

\rubric{On Shavuot:}
חַג הַשָּׁבוּעוֹת הַזֶּה, בְּיוֹם טוֹב מִקְרָא קֹדֶשׁ הַזֶּה,

\rubric{On Rosh HaShanah:}
הַזִּכָּרוֹן הַזֶּה, בְּיוֹם טוֹב מִקְרָא קֹדֶשׁ הַזֶּה,

\rubric{On Sukkot:}
חַג הַסֻּכּוֹת הַזֶּה, בְּיוֹם [טוֹב] מִקְרָא קֹדֶשׁ הַזֶּה,

\rubric{On Shemini Atzeret:}
שְׁמִינִי חַג עֲצֶרֶת הַזֶּה, בְּיוֹם טוֹב מִקְרָא קֹדֶשׁ הַזֶּה,

\rubric{On Pesah and Sukkot, say the word in brackets only on Yom Tov. Continue:}
לְרַחֵם בּוֹ עָלֵינוּ וּלְהוֹשִׁיעֵנוּ. זָכְרֵנוּ \Shem{} אֱלֹהֵינוּ בּוֹ לְטוֹבָה, וּפָקְדֵנוּ בוֹ לִבְרָכָה, וְהוֹשִׁיעֵנוּ בוֹ לְחַיִּים טוֹבִים, בִּדְבַר יְשׁוּעָה וְרַחֲמִים. חוּס וְחָנֵּנוּ, וַחֲמֹל וְרַחֵם עָלֵינוּ, וְהוֹשִׁיעֵנוּ כִּי אֵלֶיךָ עֵינֵינוּ, כִּי אֵל מֶלֶךְ חַנּוּן וְרַחוּם אָתָּה.
\end{seasonal}

\textbf{וְתִבְנֶה} יְרוּשָׁלַיִם עִירָךְ בִּמְהֵרָה בְיָמֵינוּ: בָּרוּךְ אַתָּה \Shem{} בּוֹנֵה בְּרַחֲמָיו בִּנְיַן יְרוּשָׁלַיִם עִיר הַקֹּדֶשׁ \whispered{אָמֵן}:

\textbf{בָּרוּךְ} אַתָּה \Shem{} אֱלֹהֵינוּ מֶלֶךְ הָעוֹלָם לָעַד, הָאֵל אָבִינוּ מַלְכֵּנוּ אַדִּירֵנוּ. בּוֹרְאֵנוּ. גּוֹאֲלֵנוּ. קְדוֹשֵׁנוּ. קְדוֹשׁ יַעֲקֹב. רוֹעֵנוּ רוֹעֵה יִשְׂרָאֵל. הַמֶּלֶךְ הַטּוֹב וְהַמֵּטִיב לַכֹּל. שֶׁבְּכָל יוֹם וָיוֹם הוּא הֵטִיב לָנוּ. הוּא מֵטִיב לָנוּ. הוּא יֵיטִיב לָנוּ. הוּא גְמָלָנוּ. הוּא גוֹמְלֵנוּ. הוּא יִגְמְלֵנוּ לָעַד חֵן וָחֶסֶד וְרַחֲמִים וְרֵיוַח וְהַצָּלָה וְכָל טוֹב:

הָרַחֲמָן הוּא יִשְׁתַּבַּח עַל כִּסֵּא כְבוֹדוֹ הָרַחֲמָן הוּא יִשְׁתַּבַּח בַּשָּׁמַיִם וּבָאָרֶץ הָרַחֲמָן הוּא יִשְׁתַּבַּח בָּנוּ לְדוֹר דּוֹרִים הָרַחֲמָן הוּא קֶרֶן לְעַמּוֹ יָרִים הָרַחֲמָן הוּא יִתְפָּאַר בָּנוּ לָנֶצַח נְצָחִים הָרַחֲמָן הוּא יְפַרְנְסֵנוּ בְּכָבוֹד וְלֹא בְבִזּוּי, בְּהֶתֵּר וְלֹא בְאִסּוּר, בְּנַחַת וְלֹא בְצָעַר, בְּרֶוַח וְלֹא בְצִמְצוּם הָרַחֲמָן הוּא יִתֵּן שָׁלוֹם בֵּינֵינוּ הָרַחֲמָן הוּא יִשְׁלַח בְּרָכָה רְוָחָה וְהַצְלָחָה בְּכָל מַעֲשֵׂה יָדֵינוּ הָרַחֲמָן הוּא יַצְלִיחַ אֶת דְּרָכֵינוּ הָרַחֲמָן הוּא יִשְׁבּוֹר עֹל גָּלוּת מְהֵרָה מֵעַל צַוָּארֵנוּ הָרַחֲמָן הוּא יוֹלִיכֵנוּ מְהֵרָה קוֹמְמִיּוּת בְּאַרְצֵנוּ הָרַחֲמָן הוּא יִרְפָּאֵנוּ רְפוּאָה שְׁלֵמָה רְפוּאַת הַנֶּפֶשׁ וּרְפוּאַת הַגּוּף הָרַחֲמָן הוּא יִפְתַּח לָנוּ אֶת יָדוֹ הָרְחָבָה הָרַחֲמָן הוּא יְבָרֵךְ כָּל אֶחָד וְאֶחָד מִמֶּנּוּ בִּשְׁמוֹ הַגָּדוֹל כְּמוֹ שֶׁנִּתְבָּרְכוּ אֲבוֹתֵינוּ אַבְרָהָם יִצְחָק וְיַעֲקֹב בַּכֹּל מִכֹּל כֹּל. כֵּן יְבָרֵךְ אוֹתָנוּ יַחַד בְּרָכָה שְׁלֵמָה. וְכֵן יְהִי רָצוֹן וְנֹאמַר אָמֵן הָרַחֲמָן הוּא יִפְרוֹשׂ עָלֵינוּ סֻכַּת שְׁלוֹמוֹ:

\begin{seasonal}{On Shabbat:}
הָרַחֲמָן הוּא יַנְחִילֵנוּ עוֹלָם שֶׁכֻּלּוֹ שַׁבָּת וּמְנוּחָה לְחַיֵּי הָעוֹלָמִים.
\end{seasonal}

\begin{seasonal}{On Rosh Hodesh:}
הָרַחֲמָן הוּא יְחַדֵּשׁ עָלֵינוּ אֶת הַחֹדֶשׁ הַזֶּה לְטוֹבָה וְלִבְרָכָה.
\end{seasonal}

\begin{seasonal}{On Rosh HaShanah:}
הָרַחֲמָן הוּא יְחַדֵּשׁ עָלֵינוּ אֶת הַשָּׁנָה הַזֹּאת לְטוֹבָה וְלִבְרָכָה.
\end{seasonal}

\begin{seasonal}{On Sukkot:}
הָרַחֲמָן הוּא יְזַכֵּנוּ לֵישֵׁב בְּסֻכַּת עוֹרוֹ שֶׁל לִוְיָתָן. הָרַחֲמָן הוּא יַשְׁפִּיעַ עָלֵינוּ שֶׁפַע קְדֻשָּׁה וְטָהֳרָה מִשִּׁבְעָה אֻשְׁפִּיזִין עִלָּאִין קַדִּישִׁין, זְכוּתָם תְּהֵא מָגֵן וְצִנָּה עָלֵינוּ.
\end{seasonal}

\begin{seasonal}{On Hol HaMoed:}
הָרַחֲמָן הוּא יַגִּיעֵנוּ לְמוֹעֲדִים וּרְגָלִים אֲחֵרִים הַבָּאִים לִקְרָאתֵנוּ לְשָׁלוֹם.
\end{seasonal}

\begin{seasonal}{On Yom Tov:}
הָרַחֲמָן הוּא יַנְחִילֵנוּ יוֹם שֶׁכֻּלּוֹ טוֹב.
\end{seasonal}

\textbf{הָרַחֲמָן} הוּא יִטַּע תּוֹרָתוֹ וְאַהֲבָתוֹ בְּלִבֵּנוּ וְתִהְיֶה יִרְאָתוֹ עַל פָּנֵינוּ לְבִלְתִּי נֶחֱטָא. וְיִהְיוּ כָל מַעֲשֵׂינוּ לְשֵׁם שָׁמָיִם:

\rubric{A guest adds this blessing for the host:}
הָרַחֲמָן הוּא יְבָרֵךְ אֶת הַשֻּׁלְחָן הַזֶּה שֶׁאָכַלְנוּ עָלָיו, וִיסַדֵּר בּוֹ כָּל מַעֲדַנֵּי עוֹלָם, וְיִהְיֶה כְּשֻׁלְחָנוֹ שֶׁל אַבְרָהָם אָבִינוּ. כָּל רָעֵב מִמֶּנּוּ יֹאכַל, וְכָל צָמֵא מִמֶּנּוּ יִשְׁתֶּה. וְאַל יֶחְסַר מִמֶּנּוּ כָּל טוּב לָעַד וּלְעוֹלְמֵי עוֹלָמִים, אָמֵן. הָרַחֲמָן הוּא יְבָרֵךְ אֶת בַּעַל הַבַּיִת הַזֶּה וּבַעַל הַסְּעוּדָה הַזֹּאת, הוּא וּבָנָיו וְאִשְׁתּוֹ וְכָל אֲשֶׁר לוֹ, בְּבָנִים שֶׁיִּחְיוּ וּבִנְכָסִים שֶׁיִּרְבּוּ. בָּרֵךְ \Shem{} חֵילוֹ וּפֹעַל יָדָיו תִּרְצֶה. וְיִהְיוּ נְכָסָיו וּנְכָסֵינוּ מֻצְלָחִים וּקְרוֹבִים לָעִיר. וְאַל יִזְדַּקֵּק לְפָנָיו וְלֹא לְפָנֵינוּ שׁוּם דְּבַר חֵטְא וְהִרְהוּר עָוֹן. שָׂשׂ וְשָׂמֵחַ כָּל הַיָּמִים בְּעֹשֶׁר וְכָבוֹד מֵעַתָּה וְעַד עוֹלָם. לֹא יֵבוֹשׁ בָּעוֹלָם הַזֶּה וְלֹא יִכָּלֵם לָעוֹלָם הַבָּא. אָמֵן, כֵּן יְהִי רָצוֹן.

\rubric{Personal requests may be added here.}
\textbf{הָרַחֲמָן} הוּא יְחַיֵּינוּ וִיזַכֵּנוּ וִיקָרְבֵנוּ לִימוֹת הַמָּשִׁיחַ וּלְבִנְיַן בֵּית הַמִּקְדָּשׁ וּלְחַיֵּי הָעוֹלָם הַבָּא.

\rubric{On Shabbat, Yom Tov, and days with Musaf, say \foreignlanguage{hebrew}{מִגְדּוֹל} instead of \foreignlanguage{hebrew}{מַגְדִּיל} at the start of this paragraph.}
\textbf{מַגְדִּיל} יְשׁוּעוֹת מַלְכּוֹ. וְעֹשֶׂה חֶסֶד לִמְשִׁיחוֹ לְדָוִד וּלְזַרְעוֹ עַד עוֹלָם: כְּפִירִים רָשׁוּ וְרָעֵבוּ. וְדוֹרְשֵׁי \Shem{} לֹא יַחְסְרוּ כָל טוֹב: נַעַר הָיִיתִי גַּם זָקַנְתִּי וְלֹא רָאִיתִי צַדִּיק נֶעֱזָב. וְזַרְעוֹ מְבַקֵּשׁ לָחֶם: כָּל הַיּוֹם חוֹנֵן וּמַלְוֶה וְזַרְעוֹ לִבְרָכָה: מַה שֶּׁאָכַלְנוּ יִהְיֶה לְשָׂבְעָה. וּמַה שֶּׁשָּׁתִינוּ יִהְיֶה לִרְפוּאָה. וּמַה שֶּׁהוֹתַרְנוּ יִהְיֶה לִבְרָכָה כְּדִכְתִיב וַיִּתֵּן לִפְנֵיהֶם וַיֹּאכְלוּ וַיּוֹתִרוּ כִּדְבַר \Shem{}: בְּרוּכִים אַתֶּם \Shem{לַ}. עוֹשֵׂה שָׁמַיִם וָאָרֶץ: בָּרוּךְ הַגֶּבֶר אֲשֶׁר יִבְטַח \Shem{בַּ}. וְהָיָה \Shem{} מִבְטָחוֹ: \Shem{} עֹז לְעַמּוֹ יִתֵּן. \Shem{} יְבָרֵךְ אֶת עַמּוֹ בַשָּׁלוֹם:

כִּי הִשְׂבִּיעַ נֶפֶשׁ שֹׁקֵקָה, וְנֶפֶשׁ רְעֵבָה מִלֵּא טוֹב: הוֹדוּ \Shem{לַ} כִּי טוֹב, כִּי לְעוֹלָם חַסְדּוֹ: הוֹדוּ \Shem{לַ} כִּי טוֹב, כִּי לְעוֹלָם חַסְדּוֹ: הַשָּׁמַיִם שָׁמַיִם \Shem{לַ}, וְהָאָרֶץ נָתַן לִבְנֵי אָדָם:

עוֹשֶׂה שָׁלוֹם בִּמְרוֹמָיו הוּא בְרַחֲמָיו יַעֲשֶׂה שָׁלוֹם עָלֵינוּ וְעַל כָּל עַמּוֹ יִשְׂרָאֵל וְאִמְרוּ אָמֵן:

\begin{textbox}{Ya Komimos}
Ya komimos i bevimos i al Dyo santo Baruḥ U uvaruḥ shemo bendishimos.
Ke mos dyo i mos dara pan para komer, i panyos para vistir, i anyos, munchos
i buenos, para bivir. El Padre el grande ke mande al chiko, asegun tenemos de
menester para muestras kazas i para muestros ijos. El Dyo mos oyga, i mos
aresponda, i mos apiade por su nombre el grande, ke somos almikas sin pekado.
\textit{Odu L’Adonai ki tov ki leolam ḥasdo. Odu L’Adonai ki tov ki leolam
ḥasdo.} Syempre mejor, nunka peor, nunka mos manke en la meza del Kriador.
Amen.
\end{textbox}
